% !TEX program = LuaLaTeX
\DocumentMetadata{lang=en,pdfstandard=ua-2,tagging=on,
  tagging-setup={math/setup=mathml-SE}}
\documentclass{article}
\usepackage[margin=1.1in]{geometry}
\usepackage{hyperref}
\usepackage{xurl}
\title{PDF accessibility}
\date{\today}
\author{Tim Prescott}

\begin{document}

\maketitle\thispagestyle{plain}

Making a \LaTeX\ document accessible:
\begin{itemize}
\item Get the latest versions of everything.  Install the 2025 version, and update your packages and executables (maybe see \url{https://tex.stackexchange.com/questions/55437/} for the updating part).
%\item As the first command (before \verb|\documentclass|), use \verb|\RequirePackage{nag}|.  Resolve everything it complains about, then delete that line.  Then continue on.
\item As the first command (before \verb|\documentclass|), use
\begin{verbatim}
\DocumentMetadata{lang=en,pdfstandard=ua-2,tagging=on,
  tagging-setup={math/setup=mathml-SE,math/alt/use}}
\end{verbatim}
If this won't be on Blackboard, remove the \verb|,math/alt/use|.
A useful key is \texttt{check-tagging-status}, which will report the tagging status of your class and packages.
See \texttt{texdoc documentmetadata-support-code} for more information.
\item Give your document a \verb|\title|
\item Any \verb|\includegraphics|, \verb|\tikz|, or \verb|{tikzpicture}| or \verb|{picture}| environments need \verb|alt={...}| in the options (if the image is truly decorative, you can use the option \texttt{artifact} instead of \texttt{alt}; a different option is \verb|actualtext={...}| if your image is really text in disguise).

``Best practices'' recommend that alt text has at most two sentences or 120 characters, and to keep in mind that alt text is generally invisible to most users (also don't start with ``image'', since screen readers often provide the word ``image'' on their own).  If you need significantly more than that limit, you're supposed to use the alt text to point to the longer description.
%  (The TTaDA courses ``Making Simple/Complex Images Accessible'' boil down to those two sentences.)
\item If you have a tabular, see \texttt{texdoc latex-lab-table} for additional commands you'll need.  See also \url{https://latex3.github.io/tagging-project/documentation/prototype-usage-instructions.html}.  Be forewarned that Ally has improperly complained about my \texttt{tagging=presentation}.
\item Lua\LaTeX\ works better than pdf\LaTeX; Xe\LaTeX\ doesn't really work.  The command \texttt{lualatex-dev} is the even more up to date version, which might help.
%\item Try to avoid low level (La)\TeX\ hacks and stick to package commands.  If you can \verb|\newcommand| something you'll be doing regularly, then you'll be able to edit that command instead of every instance where it was used.
%\item \verb|\normalsize| UND green (RGB 0,158,68\textasciitilde Hsb 146,1,0.62) and orange (RGB 255,103,31\textasciitilde Hsb 19,0.88,1) do not sufficiently contrast with white.  If you add 13\% and 22\% black (respectively), you get RGB 0,137,59\textasciitilde Hsb 146,1,0.53 and RGB 198,80,24\textasciitilde Hsb 19,0.88,0.88, which may sufficiently contrast.
\end{itemize}

\noindent
\verb|\DocumentMetadata| greatly slows down the compilation process with \texttt{pgfplots} by a factor of 7 or so.
% (compilations for the Calculus book have gone from minutes to hours).
You can try changing the first command to use \texttt{tagging=off}.
%I've also gone to compiling the Calculus book on UND's high performance computer, so that I can start the process and come back later.

General class and package status is available at \url{https://latex3.github.io/tagging-project/tagging-status/} (this is what gets reported by the \texttt{check-tagging-status} key), see also \url{https://latex3.github.io/tagging-project/documentation/prototype-usage-instructions.html}.

Important package/class substitutions:
\begin{itemize}
\item \texttt{enumitem} and \texttt{enumerate} need to change to \texttt{enumext} (but the default versions now accept optional arguments so you may no longer need a package).
\item \texttt{marginnote}, \texttt{marginfit}, and \texttt{marginfix} need to become \texttt{marginalia} or \verb|\marginpar|.
\item \texttt{beamer} needs to change to \texttt{ltx-talk}; documentation at \texttt{texdoc ltx-talk}.
\end{itemize}

UND is using Ally for Blackboard and Siteimprove (\url{https://www.siteimprove.com/platform/}) for the rest of the \texttt{UND.edu} website.  These don't appear too difficult to pass.  If you want to do better, you can use the much stricter VeraPDF (\url{https://verapdf.org/software/}).
% and PDF Accessibility Checker (aka PAC, \url{https://pac.pdf-accessibility.org/en}).  The latter cannot validate UA-2/2.0 (so you'll need to remove the \texttt{math/setup=mathml-SE} and keep the \texttt{math/alt/use}).

%I have accessibilized the Track Meet tests for the last two contests, and used \LaTeX ML to create html versions.
%I have also accessibilized the Calculus book and Cherney's Linear Algebra, with the following items to note:
%\begin{itemize}
%\item I needed a low level work around for media9's \verb|\includemedia|
%\item To please PAC, I needed a low level work around for tables
%\item PAC still complains about irregular tables (those using \verb|\multicolumn| or \verb|\multirow|).  It's possible to work around this (see \url{https://github.com/latex3/tagging-project/discussions/930})
%\item PAC complains about pgfplots (see \url{https://github.com/pgf-tikz/pgfplots/issues/507})
%\end{itemize}

Other resources:
\begin{description}
\item[PAC:] \url{https://pac.pdf-accessibility.org/en} is an alternative to VeraPDF.  It cannot validate UA-2/2.0 (so you'll need to remove the \texttt{math/setup=mathml-SE} and keep the \texttt{math/alt/use}).  It also needs a low level work around for tables, and still complains about irregular tables (those using \verb|\multicolumn| or \verb|\multirow|).  It's possible to work around this (see \url{https://github.com/latex3/tagging-project/discussions/930}) and a bug in pgfplots (see \url{https://github.com/pgf-tikz/pgfplots/issues/507})
\item[PDFix:] \url{https://pdfix.net/} is much stricter than Ally.  It mostly uses VeraPDF behind the scenes.
\item[\LaTeX ML:] \url{https://math.nist.gov/~BMiller/LaTeXML/} converts \TeX\ to html.  This is what ArXiv uses to create internet versions.  The result is more accessible, but the process is more difficult.
\item[PreTeXt:] \url{https://pretextbook.org/} converts its own xml syntax to html.  You need to learn a new language.
\item[Pandoc:] \url{https://pandoc.org/} converts to and from various formats.
\item[General information:] \url{https://www.accessiblemath.info/accessible_math_documentation_page.html}
\item[This document:] \url{https://sites.und.edu/timothy.prescott/accessible}
\item[Before and after example:] \url{https://www.youtube.com/watch?v=Eu_qM53tInw&t=5900s} is Frank Mittelbach's talk from the June '25 \TeX\ conference. 
\end{description}

If you have any questions, please let me know.

\end{document}
